\section{全链条关系表：存在、关系强度与估值资格是三件不同的事}

本报告梳理了 21 个产业链节点和 9 条公司关系。它们覆盖锂精矿资源、锂盐加工、民爆产品、爆破服务、运输和下游需求，但只有历史上已形成收入和毛利的锂产品、民爆产品、爆破服务可进入基准盈利桥。资源、客户和项目关系不是无效信息；它们是需要进一步披露才能被量化的观察项。

\begin{exhibitbox}[表 A-7：关键资源与产业关系的证据矩阵]
\begin{tabularx}{\exhibitboxwidth}{L{2.7cm}L{2.5cm}L{2.5cm}X}
\toprule
对象 & 已验证关系 & 未验证的关键变量 & 对模型的处理 \\
\midrule
李家沟 / 能投锂业 & 37.25\% 权益及优先供应安排 & 实际交付量、成本、自供比例、期限 & 零额外估值信用 \\
KMC & 设计采选能力有披露 & 实际供应、成本、权益与现金贡献 & 零额外估值信用 \\
锂盐加工 & 近 13 万吨设计口径、2025A 销量 6.95 万吨 & 当期利用率、产品结构、实际实现价 & 仅以历史销量+House 情景建模 \\
民爆产品 & 2025A 许可产能 26.25 万吨、利用率 83.74\%；2026 年 4 月许可更新 33 万吨 & 新许可对应实际利用率、量价、回款 & 不将新产能乘旧利用率 \\
爆破服务 & 2025A 收入与毛利已实现 & 项目订单、进度、回款与分部净利 & 温和增长，现金单列约束 \\
运输 & 2025A 收入已披露 & 独立毛利、客户结构与现金贡献 & 使用派生毛利率代理，低权重 \\
\bottomrule
\end{tabularx}
\end{exhibitbox}

\section{客户链审计：不把名称或主题变成订单}

客户、下游电池厂、矿山项目与政策主题常被用于强化锂盐故事。本案的客户链审计共记录 10 个观察项，但没有足以将某客户的收入、价格、期限、订单量和回款条款放入前瞻模型的公司级披露。因此，任何客户或项目名称都不产生基准 EPS；正式披露若满足金额、期限和履约三项之一以上，才可从观察项转入模型。

\begin{exhibitbox}[表 A-8：证据资格与模型权限]
\begin{tabularx}{\exhibitboxwidth}{L{3.0cm}X X}
\toprule
证据级别 & 典型信息 & 模型权限 \\
\midrule
一级：公司已实现并经审计 & 历史分部收入、成本、毛利、销量、现金流 & 可作为基准起点 \\
二级：公司明确披露但未实现 & 权益、许可、设计能力、优先供应安排 & 只作战略与上行情景观察 \\
三级：外部行业或同业资料 & 商品价格、下游销量、同行估值 & 只作方向和交叉校验 \\
四级：名称、传闻或不可复算线索 & 客户、项目、资源成本、订单讨论 & 不进入收入、利润或目标价 \\
\bottomrule
\end{tabularx}
\end{exhibitbox}

\section{监管与运营提醒}

民爆业务的许可容量在 2026 年 4 月更新为 33 万吨，而 2025 年年报披露的 26.25 万吨许可容量和 83.74\% 利用率对应的是历史年度。把更新后的容量直接乘以前一年的利用率，会夸大实际产出；同理，设计锂盐能力不能替代实际出货。研究员应在 H1 后重新核对许可证、产量、收入与现金四者是否同口径。

\sourcenote{产业链地图、关系表、客户链审计以及公司年报与公告。}
